%% Chapter 4 — OCEL Conformance as a Measure-Theoretic Framework

\chapter{OCEL Conformance as a Measure-Theoretic Framework}
\label{chap:ocel}

\epigraph{%
  ``If you cannot observe the process, you cannot manage the process.''
}{--- Wil van der Aalst, \textit{Process Mining}, 2016}

Chapter~\ref{chap:oracle} established that the 483-test suite is an
unfakeable oracle: no algorithmically simpler function can pass it.  But
test-suite passage is a static guarantee---it certifies behaviour on a fixed
corpus of inputs.  Runtime executions generate \emph{traces}, and a breed
that passes every laboratory test could, in principle, follow an unlawful
execution path during production.  This chapter closes that gap by
constructing a measure-theoretic framework around the Object-Centric Event
Log (OCEL) \cite{ocel2standard,kusters2025ocpq} emitted by every \textsc{wasm4pm} execution, proving a
conformance certificate theorem that links static correctness to dynamic
fitness, and establishing that the Breed Validation Certificate (BVC)
machinery enforces \textbf{L0 wrapper fitness} $\varphi_{\mathrm{L0}} = 1.0$ as an
admission gate for every production execution. L1 breed-inference alignment fitness
(token replay against the declared breed OCPN) has been measured as step-coverage
$= 1.0$ for MYCIN, Prolog, and STRIPS; full alignment fitness via pm4py replay
is scheduled for proof-pack v2 and currently carries \textsc{partial-alive} standing
(see \texttt{ocel/reports/\{mycin,prolog,strips\}.json}).

% ─────────────────────────────────────────────────────────────────────────────
\section{The Breed Lifecycle Object-Centric Petri Net}
\label{sec:breed-ocpn}
% ─────────────────────────────────────────────────────────────────────────────

Every execution of a breed $b_i$ passes through a fixed sequence of
lifecycle states.  We formalise this sequence as an Object-Centric Petri
Net (OCPN) and show that the OCEL emitted by
\texttt{crates/wasm4pm-cognition/src/ocel.rs} is the trace language of
this net.

\begin{definition}[Breed Lifecycle OCPN]
\label{def:breed-ocpn}
The \emph{breed lifecycle OCPN} is the tuple
\[
  \mathcal{N}_{\mathrm{breed}}
  \;=\;
  \bigl(P,\, T,\, F,\, \lambda,\, \mu_0\bigr),
\]
where:
\begin{itemize}
  \item $P = \{p_{\mathrm{idle}},\, p_{\mathrm{dispatched}},\,
    p_{\mathrm{executing}},\, p_{\mathrm{verifying}},\,
    p_{\mathrm{committed}},\, p_{\mathrm{failed}}\}$
    is the set of \emph{places}, each representing a lifecycle state of a
    breed execution object;
  \item $T = \{t_{\mathrm{dispatch}},\, t_{\mathrm{execute}},\,
    t_{\mathrm{verify}},\, t_{\mathrm{commit}},\, t_{\mathrm{fail}}\}$
    is the set of \emph{transitions}, each labelled by the corresponding
    OCEL activity name;
  \item $F \subseteq (P \times T) \cup (T \times P)$ is the standard
    flow relation encoding the ordering
    $p_{\mathrm{idle}} \to t_{\mathrm{dispatch}} \to
     p_{\mathrm{dispatched}} \to t_{\mathrm{execute}} \to
     p_{\mathrm{executing}} \to t_{\mathrm{verify}} \to
     p_{\mathrm{verifying}} \to t_{\mathrm{commit}} \to
     p_{\mathrm{committed}}$,
    with the exception arc $t_{\mathrm{fail}}$ reachable from
    $p_{\mathrm{executing}}$ and $p_{\mathrm{verifying}}$ leading to
    $p_{\mathrm{failed}}$;
  \item $\lambda : T \to \Sigma$ maps each transition to its OCEL activity
    label in the controlled vocabulary
    $\{\texttt{breed:dispatch},\, \texttt{breed:execute},\,
      \texttt{breed:verify},\, \texttt{breed:commit},\,
      \texttt{breed:fail}\}$; and
  \item $\mu_0 = [p_{\mathrm{idle}}]$ is the initial marking (one token
    in $p_{\mathrm{idle}}$, all others empty).
\end{itemize}

Timestamps in the emitted OCEL are \emph{logical}, not wall-clock: the
epoch anchor is the constant \texttt{1970-01-01T00:00:00Z} and each event
carries a strictly increasing integer attribute \texttt{logical\_step}.
This design ensures that OCEL conformance checks are deterministic and
reproducible across machines and time zones.
\end{definition}

The implementation in \texttt{ocel.rs} materialises
$\mathcal{N}_{\mathrm{breed}}$ by emitting one OCEL event per transition
firing, recording the object identifier (the \texttt{run\_id}), the
activity label, the logical step, and the output hash at the
\texttt{breed:commit} transition.  Because \texttt{run\_id} is unique per
execution, each trace is an independent object lifecycle; the accumulation
of traces across executions forms the object-centric corpus $\mathcal{L}$
analysed in Section~\ref{sec:discovery-convergence}.

% ─────────────────────────────────────────────────────────────────────────────
\section{The OCEL Fitness Measure}
\label{sec:ocel-fitness}
% ─────────────────────────────────────────────────────────────────────────────

We now define alignment-based fitness in the object-centric setting.
Alignment-based fitness is the standard measure used in process mining
\citep{adriansyah2014aligning}: it penalises every move-on-log (event in
the trace not enabled in the model) and every move-on-model (model
transition fired without a corresponding event) by assigning unit cost to
each deviation.

\begin{definition}[Alignment-Based OCEL Fitness]
\label{def:ocel-fitness}
Let $\sigma = \langle e_1, \ldots, e_n \rangle$ be an object trace (the
sequence of OCEL events associated with a single \texttt{run\_id} object).
An \emph{alignment} of $\sigma$ against $\mathcal{N}_{\mathrm{breed}}$ is
a sequence of \emph{moves}
\[
  \gamma \;=\; \langle m_1, \ldots, m_k \rangle,
  \qquad
  m_j \in
    \bigl(\Sigma \times \{\gg\}\bigr)
    \cup
    \bigl(\{\gg\} \times T\bigr)
    \cup
    \bigl(\Sigma \times T\bigr),
\]
where $(\sigma_j, \gg)$ is a \emph{log move} (cost 1), $(\gg, t_j)$ is a
\emph{model move} (cost 1), and $(\sigma_j, t_j)$ is a \emph{synchronous
move} (cost 0), subject to the constraint that the synchronous and model
moves form a valid firing sequence from $\mu_0$ to a final marking, and the
log projection of $\gamma$ equals $\sigma$.

The \emph{fitness} of $\sigma$ is
\[
  \varphi(\sigma)
  \;=\;
  1 - \frac{\mathrm{cost}(\gamma^*(\sigma))}{|\sigma| + |T_{\mathrm{model}}|},
\]
where $\gamma^*(\sigma)$ is the minimum-cost alignment of $\sigma$ and
$T_{\mathrm{model}}$ is the minimal sequence of model-side transitions in
any complete firing sequence.  The \emph{corpus fitness} is
\[
  \Phi(\mathcal{L})
  \;=\;
  \frac{1}{|\mathcal{L}|} \sum_{\sigma \in \mathcal{L}} \varphi(\sigma).
\]
\end{definition}

\begin{axiom}[BVC Fitness Requirement]
\label{ax:bvc-fitness}
Every production execution of a breed $b_i$ must satisfy
$\varphi(\sigma) = 1.0$; that is, the alignment of its object trace
against $\mathcal{N}_{\mathrm{breed}}$ must be the identity (no log moves,
no model moves, only synchronous moves).  An execution that does not satisfy
this axiom is \emph{non-conforming} and must be refused.
\end{axiom}

Axiom~\ref{ax:bvc-fitness} is not a statistical target but a hard
admission criterion.  Its enforcement mechanism is the \texttt{run\_breed}
choke point described in the next section.

% ─────────────────────────────────────────────────────────────────────────────
\section{The Conformance Certificate Theorem}
\label{sec:conformance-certificate}
% ─────────────────────────────────────────────────────────────────────────────

\begin{theorem}[Conformance Certificate]
\label{thm:conformance-certificate}
Let $\mathcal{N}_{\mathrm{breed}}$ be as in Definition~\ref{def:breed-ocpn}
and let $\mathcal{L}$ be the OCEL corpus emitted by \textsc{wasm4pm}.
Then:

\begin{enumerate}[label=(\roman*)]
  \item \textbf{Soundness.}  Every trace $\sigma \in \mathcal{L}$
    satisfies $\varphi(\sigma) = 1.0$ if and only if the execution that
    generated $\sigma$ completed the full dispatch--execute--verify--commit
    path without deviation.

  \item \textbf{Completeness.}  If an execution completes the full path,
    the OCEL emitter in \texttt{ocel.rs} produces exactly one event per
    transition in $\mathcal{N}_{\mathrm{breed}}$, so the resulting trace
    $\sigma$ has a trivial (zero-cost) alignment.

  \item \textbf{Fitness Universality.}  Non-conforming traces---those with
    $\varphi(\sigma) < 1.0$---are refused at the \texttt{run\_breed} choke
    point implemented at lines 70--77 of
    \texttt{crates/wasm4pm-cognition/src/wasm.rs}.  Consequently,
    $\Phi(\mathcal{L}) = 1.0$ for all production corpora.
\end{enumerate}
\end{theorem}

\begin{proof}
\textbf{Clause (i) --- Soundness.}
Fitness $\varphi(\sigma) = 1.0$ implies $\mathrm{cost}(\gamma^*) = 0$,
which in turn implies no log moves and no model moves in the optimal
alignment.  A zero-cost alignment is a purely synchronous sequence; by the
definition of $\mathcal{N}_{\mathrm{breed}}$, the unique synchronous
firing sequence from $\mu_0$ to the final marking visits the transitions in
the order $t_{\mathrm{dispatch}},\, t_{\mathrm{execute}},\,
t_{\mathrm{verify}},\, t_{\mathrm{commit}}$.  Conversely, if the execution
followed this path, the OCEL emitter (clause (ii)) produces exactly the
events that label these transitions, so the alignment is the identity map
and the cost is zero.  \qquad $\square_{(i)}$

\medskip\noindent
\textbf{Clause (ii) --- Completeness.}
Inspect the control flow of \texttt{ocel.rs}: each of the five OCEL
activities is emitted precisely once per execution object (the
\texttt{run\_id} is allocated at dispatch, and subsequent event-emit calls
reference the same identifier).  The logical step counter is incremented
atomically before each emit, guaranteeing strict monotonicity.  A trace
produced by this emitter therefore contains exactly $|T| = 5$ events (or
4 if the fail transition fires instead of commit) in the order mandated by
the flow relation $F$.  The minimum-cost alignment of such a trace against
$\mathcal{N}_{\mathrm{breed}}$ is the identity, giving cost 0 and fitness
1.0.  \qquad $\square_{(ii)}$

\medskip\noindent
\textbf{Clause (iii) --- Fitness Universality.}
The \texttt{run\_breed} function at lines 70--77 of \texttt{wasm.rs}
evaluates the conformance check synchronously before returning the
\texttt{ContractResult} to the caller.  If the check determines that the
pending trace would have $\varphi < 1.0$---for example because a required
transition was skipped or an unexpected activity was recorded---the function
returns an error with status \texttt{conformance\_rejected} and does not
advance the breed execution to the commit state.  The OCEL corpus
$\mathcal{L}$ therefore contains only traces that passed this gate.  Every
trace in $\mathcal{L}$ has $\varphi(\sigma) = 1.0$ by clause (i), so
\[
  \Phi(\mathcal{L})
  = \frac{1}{|\mathcal{L}|} \sum_{\sigma \in \mathcal{L}} 1.0
  = 1.0.
\]
\qquad $\square_{(iii)}$
\end{proof}

% ─────────────────────────────────────────────────────────────────────────────
\section{Process Mining Over the Accumulated Corpus}
\label{sec:discovery-convergence}
% ─────────────────────────────────────────────────────────────────────────────

Theorem~\ref{thm:conformance-certificate} guarantees that every trace in
$\mathcal{L}$ is individually conformant.  A complementary question is
whether the \emph{discovered} process model---obtained by applying a
discovery algorithm to $\mathcal{L}$---converges to
$\mathcal{N}_{\mathrm{breed}}$ as the corpus grows.

\begin{proposition}[Discovery Convergence]
\label{prop:discovery-convergence}
Let $\mathcal{L}_n$ denote the OCEL corpus after $n$ production executions,
all passing the \texttt{run\_breed} gate, and let $\hat{\mathcal{N}}_n$ be
the Petri net discovered from $\mathcal{L}_n$ by the inductive miner
\citep{leemans2013discovering}.  Then:
\[
  \lim_{n \to \infty} \fitness\bigl(\hat{\mathcal{N}}_n,\, \mathcal{L}_n\bigr)
  \;=\; 1.0,
\]
and for sufficiently large $n$, $\hat{\mathcal{N}}_n$ is isomorphic to
$\mathcal{N}_{\mathrm{breed}}$ as a labelled Petri net.
\end{proposition}

\begin{proof}
By Theorem~\ref{thm:conformance-certificate}(iii), every trace in
$\mathcal{L}_n$ follows the path
$t_{\mathrm{dispatch}} \to t_{\mathrm{execute}} \to t_{\mathrm{verify}}
\to t_{\mathrm{commit}}$ (with the fail path as the only exception arc).
The inductive miner constructs a process tree by identifying the most
frequent directly-follows relations; because all conformant traces share
the same directly-follows structure (the linear chain is the only path to
commit), the directly-follows graph over $\mathcal{L}_n$ converges to the
single-path graph of $\mathcal{N}_{\mathrm{breed}}$ as $n \to \infty$.
Conversion of the converged process tree back to a Petri net via the
standard tree-to-net transformation \citep{leemans2013discovering} yields a
net isomorphic to $\mathcal{N}_{\mathrm{breed}}$.  Fitness of the
discovered net against the corpus from which it was discovered approaches
1.0 by standard inductive-miner guarantees under the assumption that the
trace language is a regular superset of the target language---which holds
here because the trace language is a singleton (one linear path to commit,
one to fail).
\end{proof}

Proposition~\ref{prop:discovery-convergence} is the van der Aalst
constitution in concrete form: the process-truth engine does not merely
verify static tests; it \emph{discovers} the process from evidence and
confirms that the discovered model matches the declared model.  Any
divergence between $\hat{\mathcal{N}}_n$ and $\mathcal{N}_{\mathrm{breed}}$
is a defect, not a discrepancy.

% ─────────────────────────────────────────────────────────────────────────────
\section{Summary}
\label{sec:ocel-summary}
% ─────────────────────────────────────────────────────────────────────────────

This chapter constructed the breed lifecycle OCPN
(Definition~\ref{def:breed-ocpn}), defined alignment-based OCEL fitness
(Definition~\ref{def:ocel-fitness}), and mandated $\varphi = 1.0$ as a
hard admission criterion (Axiom~\ref{ax:bvc-fitness}).  The Conformance
Certificate Theorem (Theorem~\ref{thm:conformance-certificate}) established
soundness, completeness, and fitness universality via the
\texttt{run\_breed} choke point in \texttt{wasm.rs}.  Proposition~\ref{prop:discovery-convergence}
showed that the discovered process model converges to the declared model as
the corpus grows, closing the loop between static specification and dynamic
evidence.  Together, Chapters~\ref{chap:oracle} and~\ref{chap:ocel}
provide the full verification stack: unfakeable tests certify algorithmic
correctness at commit time, and OCEL conformance certifies execution
correctness at runtime.  Chapter~\ref{chap:bvc} constructs the Breed
Validation Certificate that aggregates both layers into a single
cryptographically bound receipt.
